% !TeX program = xelatex

\documentclass[]{article}
\usepackage{ragged2e}
\usepackage[onehalfspacing]{setspace}
\usepackage[a4paper, margin=0.7in]{geometry}
\usepackage{listings}
\usepackage{lstautogobble}
\usepackage{float, graphicx, subcaption}
\usepackage{xcolor} % 1. Load the color package

\usepackage[fontsize=14pt]{scrextend}

\usepackage{fontspec}
\setmainfont{Times New Roman}
\setmonofont{Iosevka Nerd Font}

\usepackage{amssymb, amsmath}
\usepackage[lite, zswash]{mtpro2}
\DeclareMathSymbol{0}{\mathalpha}{operators}{`0}
\DeclareMathSymbol{1}{\mathalpha}{operators}{`1}
\DeclareMathSymbol{2}{\mathalpha}{operators}{`2}
\DeclareMathSymbol{3}{\mathalpha}{operators}{`3}
\DeclareMathSymbol{4}{\mathalpha}{operators}{`4}
\DeclareMathSymbol{5}{\mathalpha}{operators}{`5}
\DeclareMathSymbol{6}{\mathalpha}{operators}{`6}
\DeclareMathSymbol{7}{\mathalpha}{operators}{`7}
\DeclareMathSymbol{8}{\mathalpha}{operators}{`8}
\DeclareMathSymbol{9}{\mathalpha}{operators}{`9}

% 2. Define your custom color palette
\definecolor{codekeyword}{RGB}{0, 102, 204}    % Deep Blue
\definecolor{codestring}{RGB}{0, 153, 76}     % Forest Green
\definecolor{codecomment}{RGB}{128, 128, 128} % Muted Gray
\definecolor{codebg}{RGB}{245, 245, 245}      % Light Gray Background

% 3. Apply colors to the listings settings
\lstset{
	basicstyle=\small\ttfamily,
	numbers=left,
	numberstyle=\ttfamily,
	autogobble=true,
	backgroundcolor=\color{codebg}, 
	keywordstyle=\color{codekeyword}\bfseries, 
	stringstyle=\color{codestring},    
	commentstyle=\color{codecomment}, 
	breaklines=true, 
	showstringspaces=false            
}

\begin{document}
	
	Let's introduce the first inkscape figure
	\begin{figure}[H]
		\input{drawing.pdf_tex}
	\end{figure}

	
	
	\begin{figure}
		\begin{subfigure}{0.45\textwidth}
			\def\svgwidth{8.5cm}
			\input{paths.pdf_tex}
			\caption{Paths in Complex Plane}
		\end{subfigure}
		\hfill
		\begin{subfigure}{0.55\textwidth}
			\def\svgwidth{10.5cm}
			\input{circle.pdf_tex}
			\caption{The difference between $|z|=r$ and $|z|=2r$}
		\end{subfigure}
	\end{figure}
	
	
	\begin{figure}[H]
		\def\svgwidth{15cm}
		\input{area.pdf_tex}
	\end{figure}
	
	\begin{figure}[H]
%		\def\svgwidth{15cm}
		\input{grid.pdf_tex}
	\end{figure}
	
%	\newpage
	Here some text before the TikZ
	\begin{figure}[H]
		\centering
		\includegraphics{ipetest.pdf}
	\end{figure}
	And here after the TikZ
	
	Let us do this BRO, the equation \(1 + 2 + 3 + \cdots+5\cdot\infty = \dfrac{-1}{12}\)
	\begin{lstlisting}[language=Octave, gobble=15, title=$\centerdot$\texttt{Matlab/Octave}]
		f = @(x) x^2 - 5;
		df = @(x) 2 * x; 
		
		x0 = 1;
		e = 1e-12;
		i = 0;
		
		while abs(f(x0)) > e
			x0 = x0 - f(x0) / df(x0);
			i = i + 1;
		end
		
		fprintf("The solution is %f in %d iterations", x0, i)
	\end{lstlisting}
	
\end{document}
